% Owner: Afzal. See docs/iml-planning/TEAM_WORK_SPLIT_V2.md section 8
\section{Conclusion}
A ticket records where an officer was as much as where a car was parked, and in Bengaluru officers almost disappear in the evening. On the prediction side, simple models are enough: a lasso on $\log(1+y)$ forecasts next-day station counts with macro MASE 0.697 against 0.967 for ``same weekday last week'', and most of that skill comes from a station's recent level. For next-week hotspots, four classical classifiers reach PR-AUC about 0.80 against 0.743 for ranking by last week's impact, and 0.36 to 0.37 against 0.248 on hexagons that were not hotspots the week before. Linear models match tree ensembles on both tasks.

These models predict where tickets will be written, which is a statement about enforcement workload and not about where violations occur; the evening is not recoverable from them. TODO(after E2): one sentence with the benchmark result on how much censoring hurts and how well inverse-propensity weighting recovers the truth. The study covers five months of one city and four folds, no congestion ground truth exists (the impact score is a designed proxy), and a zero count cannot distinguish ``no violation'' from ``no officer''; the intervals above capture variation over stations and hexagons but not over weeks. Natural next steps are the exposure-aware targets of the benchmark applied to real data, spatial-scale comparisons as in \citet{gao2019parking}, and a longer record that includes a full year.
